% ============================================================

\newpage

\section{The library implementation protocol}

% ============================================================

The central design principle guiding the extension was \emph{functional equivalence}: for every corruption method offered by PuckTrick, the Spark implementation must accept the same parametrisation and produce statistically equivalent output distributions, so that experiments conducted on a single-machine Pandas backend yield results that are directly comparable with those obtained in a distributed Spark cluster.  This requirement poses non-trivial engineering challenges, because Pandas and Spark embody fundamentally different computation models, and many of PuckTrick's corruption algorithms rely on random sampling, in-place mutation, and index-based row selection: operations that do not have direct analogues in the Spark API.

The implementation involved:

\begin{enumerate}

    \item Creating documentation of all existing PuckTrick corruption methods to

          catalogue their Pandas-specific dependencies.

    \item The design of a \emph{dual-backend architecture} in which each

          corruption module detects the type of the input DataFrame (Pandas

          \texttt{pd.DataFrame} or PySpark \texttt{pyspark.sql.DataFrame})

          and dispatches to the appropriate implementation.

    \item The re-implementation of each corruption method using PySpark's

          SQL-based DataFrame API, paying particular attention to

          window functions to simulate

          index-based operations.

    \item An extensive validation campaign comparing the outputs of the Pandas

          and Spark backends across all corruption types and all supported data

          types (continuous, discrete, binary, categorical string,

          categorical integer),

    \item Performance benchmarking on a multi-node Spark cluster to evaluate

          the distributed scalability of the extended library.

\end{enumerate}

